% Preamble for mathematics/ (stage-01 skeleton). Copied from plan/staged/preamble.tex (S01-M0);
% the macro set is a superset of SPLIT.md "Macros" and of the status markers used in stage 01.
% Edit here for mathematics/ only; plan/staged/preamble.tex stays the planning-side source.
\usepackage[T1]{fontenc}
\usepackage[utf8]{inputenc}
\usepackage{lmodern}
\usepackage[a4paper,margin=2.5cm]{geometry}
\usepackage{amsmath,amssymb,amsthm,mathtools}
\usepackage{booktabs,array,longtable}
\usepackage{enumitem}
\usepackage{xcolor}
\usepackage[backend=biber,style=alphabetic,maxbibnames=6]{biblatex}
\addbibresource{references.bib}
\usepackage[colorlinks=true,linkcolor=blue!50!black,citecolor=green!40!black,urlcolor=blue!50!black]{hyperref}
\setlist{itemsep=2pt,topsep=3pt}
\setlength{\emergencystretch}{3em}

%---------------------------------------------------------------- theorem environments
\theoremstyle{plain}
\newtheorem{theorem}{Theorem}[section]
\newtheorem{proposition}[theorem]{Proposition}
\newtheorem{lemma}[theorem]{Lemma}
\newtheorem{corollary}[theorem]{Corollary}
\newtheorem{conjecture}[theorem]{Conjecture}
\theoremstyle{definition}
\newtheorem{definition}[theorem]{Definition}
\newtheorem{hypothesis}[theorem]{Hypothesis}
\newtheorem{interface}[theorem]{Interface}
\newtheorem{example}[theorem]{Example}
\newtheorem{testcase}[theorem]{Test case}
\theoremstyle{remark}
\newtheorem{remark}[theorem]{Remark}

%---------------------------------------------------------------- status markers
% Usage: \status{cited}, \status{to prove}, \status{conjectural}, \status{known}, ...
\newcommand{\status}[1]{\textup{\textsf{[#1]}}}
\newcommand{\stknown}{\status{known}}          % standard / textbook, reference given
\newcommand{\stcited}{\status{cited}}          % taken from literature, precise pointer given
\newcommand{\sttoprove}{\status{to prove}}     % we must write the proof
\newcommand{\stproved}{\status{proved}}        % proof written in this plan / seed
\newcommand{\stpmc}{\status{proved modulo cited}}
\newcommand{\stsketch}{\status{sketched}}
\newcommand{\stexp}{\status{expected}}
\newcommand{\stconj}{\status{conjectural}}
\newcommand{\sthyp}{\status{hypothesis}}
\newcommand{\stopen}{\status{open}}
\newcommand{\stcomp}{\status{computed}}
\newcommand{\stformal}{\status{formalized}}    % exists in Lean (give declaration name)
% literature pointer: \lit{file-stem}{Thm 4.5, p.~12}

\newcommand{\litund}{\char`\_\penalty0\hskip0pt\relax}
\newcommand{\lit}[2]{\textup{(\pathtt{literature/#1}, #2)}}

%---------------------------------------------------------------- number systems
\newcommand{\Q}{\mathbb{Q}}\newcommand{\Z}{\mathbb{Z}}\newcommand{\R}{\mathbb{R}}\newcommand{\C}{\mathbb{C}}
\newcommand{\F}{\mathbb{F}}\newcommand{\A}{\mathbb{A}}
\newcommand{\Zp}{\mathbb{Z}_p}\newcommand{\Qp}{\mathbb{Q}_p}\newcommand{\Zq}{\mathbb{Z}_q}\newcommand{\Qq}{\mathbb{Q}_q}
\newcommand{\Zhat}{\hat{\mathbb{Z}}}\newcommand{\Qhat}{\hat{\mathbb{Q}}}
\newcommand{\cO}{\mathcal{O}}\newcommand{\cOE}{\mathcal{O}_E}

%---------------------------------------------------------------- groups, algebras
\newcommand{\GL}{\mathrm{GL}}\newcommand{\SL}{\mathrm{SL}}\newcommand{\PGL}{\mathrm{PGL}}
\newcommand{\Nrd}{\mathrm{Nrd}}\newcommand{\Trd}{\mathrm{Trd}}\newcommand{\Sym}{\mathrm{Sym}}
\newcommand{\JL}{\mathrm{JL}}\newcommand{\ord}{\mathrm{ord}}\newcommand{\Iw}{\mathrm{Iw}}
\newcommand{\Hb}{\mathfrak{H}}
\newcommand{\Bd}{B_D}                       % definite quaternion algebra over Q of discriminant D
\newcommand{\Bhat}{\hat{B}}                 % B tensor finite adeles
\newcommand{\XB}[1]{X_B(#1)}                % class set B^x \ \hat B^x / U
\newcommand{\XBq}[1]{X_{B_q}(#1)}
\newcommand{\Uzero}[1]{U_0(#1)}\newcommand{\Uone}[1]{U_1(#1)}
\newcommand{\Brandt}[2]{M(#1;#2)}           % Brandt module M(U;V)
\newcommand{\Vk}{V_k}                       % Sym^{k-2}(Z_p^2)
\newcommand{\GU}{\Gamma_U}                  % norm-one S-arithmetic group O_U[1/q]^1

%---------------------------------------------------------------- trees and graphs
\newcommand{\cT}{\mathcal{T}}               % Bruhat--Tits tree T_q
\newcommand{\G}{\mathbb{G}}                 % weighted graph / quotient graph of groups
\newcommand{\GDR}{\G^{\mathrm{DR}}}
\newcommand{\Phicomb}{\Phi^{\mathrm{comb}}}
\newcommand{\Xcomb}{X^{\mathrm{comb}}}
\newcommand{\cX}{\mathcal{X}}               % tree model of indefinite stage
\newcommand{\quot}{\backslash\backslash}    % quotient graph of groups

%---------------------------------------------------------------- Hecke, Hida
\newcommand{\T}{\mathbb{T}}\newcommand{\m}{\mathfrak{m}}
\newcommand{\Lam}{\Lambda}
\newcommand{\cM}{\mathcal{M}}
\newcommand{\Mord}{\cM^{\mathrm{ord}}_B}
\newcommand{\Xord}{\cX^{\mathrm{ord}}}
\newcommand{\Tord}{\mathcal{T}^{\mathrm{ord}}} % Hida's ordinary Tate module
\newcommand{\fa}{\mathfrak{a}}\newcommand{\fb}{\mathfrak{b}}
\newcommand{\cc}{\mathfrak{c}}\newcommand{\tq}{\mathfrak{t}}

%---------------------------------------------------------------- hypotheses
\newcommand{\HypX}{\mathrm{Hyp.X}}
\newcommand{\HypXchar}{\mathrm{Hyp.X.char}}
\newcommand{\HypXgal}{\mathrm{Hyp.X.gal}}
\newcommand{\HypXiso}{\mathrm{Hyp.X}^{\mathrm{iso}}}

%---------------------------------------------------------------- operators
\DeclareMathOperator{\Hom}{Hom}\DeclareMathOperator{\coker}{coker}\DeclareMathOperator{\End}{End}
\DeclareMathOperator{\Ann}{Ann}\DeclareMathOperator{\rk}{rk}\DeclareMathOperator{\Stab}{Stab}
\DeclareMathOperator{\tr}{tr}\DeclareMathOperator{\Jac}{Jac}\DeclareMathOperator{\Fitt}{Fitt}

%---------------------------------------------------------------- macros shared by several stages
% (moved here from the stage files; the stages' \providecommand lines are then no-ops)
\providecommand{\Tpath}[1]{\pathtt{#1}}
\providecommand{\BG}{\mathrm{BG}}          % bipartite Brandt graph of (U,q)
\providecommand{\Xone}{X^1_B}              % norm-one class set
\providecommand{\HypSA}{\mathrm{Hyp.SA}}   % strong approximation hypothesis
\providecommand{\HypHS}{\mathrm{Hyp.HS}}   % Hasse--Schilling norm hypothesis
\providecommand{\NUp}{N_{U_p}}             % the U_p normalisation fixed in stage 01
\providecommand{\Kp}[1]{K_{#1}}            % K_s = U_1(p^s)_p
\providecommand{\Kpp}[1]{K'_{#1}}          % K'_s = U_1(p^s)_p \cap U_0(p^{s+1})_p
\providecommand{\TordB}{\T^{\mathrm{ord}}_B}   % ordinary definite Hecke algebra (NOT \Tord)
\providecommand{\cH}{\mathcal{H}}          % abstract Hecke algebra H_s of Hyp.X
\providecommand{\Xq}{X_q}                  % character group at q
\providecommand{\Ros}{\mathsf{Ros}}          % Rosati dual (distinct from \vee = dual variety/morphism)
\providecommand{\Frob}{\mathrm{Frob}}
\providecommand{\St}{\mathrm{St}}          % Steinberg representation

%---------------------------------------------------------------- provisional changes (global Fable seam audit)
% \provisional{Fk/Dn: what changed; before: ...} marks a change made provisionally after
% audits/global-fable-seam.md, pending Fabian's review (see audits/global-fable-resolution.md).
% Revert a change by restoring the "before" text and deleting the note.
\providecommand{\provisional}[1]{\textcolor{orange!75!black}{\textsf{\small[provisional: #1]}}}

%---------------------------------------------------------------- Lean and status markers (mathematics/ only)
% \lean{Langlandscape.Name}: cross-reference to a formalized declaration (typeset in small \texttt).
% \pathtt{...}: typewriter path or Lean name with line-break points after /, ., \_ (round R5; no external files).
\ExplSyntaxOn
\tl_new:N \l__llpath_tl
\cs_new_protected:Npn \__llpath_typeset:n #1
  {
    \tl_set:Nn \l__llpath_tl {#1}
    \tl_replace_all:Nnn \l__llpath_tl { / } { /\penalty\exhyphenpenalty\hskip0pt\relax }
    \tl_replace_all:Nnn \l__llpath_tl { . } { .\penalty\exhyphenpenalty\hskip0pt\relax }
    \tl_replace_all:Nnn \l__llpath_tl { \_ } { \_\penalty\exhyphenpenalty\hskip0pt\relax }
    \tl_replace_all:Nnn \l__llpath_tl { - } { -\penalty\exhyphenpenalty\hskip0pt\relax }
    \tl_replace_all:Nnn \l__llpath_tl { : } { :\penalty\exhyphenpenalty\hskip0pt\relax }
    \texttt { \tl_use:N \l__llpath_tl }
  }
\NewDocumentCommand \pathtt { m } { \__llpath_typeset:n {#1} }
\ExplSyntaxOff
\newcommand{\lean}[1]{\textup{\small(Lean: \pathtt{#1})}}
% \leanplanned: the result is planned for Lean but not yet formalized.
\newcommand{\leanplanned}{\textup{\small(Lean: planned)}}
% \stformalized: status marker, exists in Lean (same meaning as \stformal; give \lean{} next to it).
\newcommand{\stformalized}{\status{formalized}}
% Loose lines rather than overfull ones (round R5).
\setlength{\emergencystretch}{3em}
% Bibliography notes contain long local paths: allow loose lines there (round R5).
\AtBeginDocument{\ifdefined\bibsetup\appto\bibsetup{\sloppy\emergencystretch=3em}\fi}
